\subsection{Situated Kinodynamic Manifold}\label{sec:edge-bundles}

Given the highly varied nature of prehensile and nonprehensile motion primitives, it may not always be possible to simulate multiple control parameters for each motion primitive within a reasonable amount of time. Moreover, doing this repeatedly for new planning problems with no memory of past experience adds to the computational cost of planning kinodynamically--feasible paths in high--dimensional workspaces.

Our work extends the concept of kinodynamic edge bundles to the multimodal domain by sampling edges based on forward simulation of the robot dynamics model \cite{shome2021asymptotically}. We term this \textsl{kinodynamic manifold} since it represents a generic set of primitive--agnostic edges that leads us to the definition of a \textsl{situated kinodynamic manifold}. More specifically, a \textsl{kinodynamic manifold} is a set of valid rollouts of the manipulator dynamics model generated by applying randomly sampled joint velocities at random joint configurations for a sampled period of time, thereby connecting random start configurations to resultant end configurations:
\begin{center}
% \begin{align}
$\mathcal{K}_{MP} = \{(x_s, x_f, u, \Delta t) \mid x_f = f_{MP}(x_s, u, \Delta t)\}$
% \end{align}
\end{center}
where $x_s, x_f \in \mathcal{X}$, $u \in \mathcal{U}_{MP}$, and $\Delta t \in \mathbb{R}$.
% such that $f_{MP}(x_s, u, \Delta t) = x_f$.
% where $|\mathcal{E}_{MP}| = m$.


% \input{definitions/kinodynamic-manifold.tex}

The manipulator dynamics model (\Cref{algo:generate-kinodynamic-manifold}, L12) is defined (here, for a 7DoF robot arm) as 
\begin{center}
$\tau = M(q)\ddot{q} + C(q, \dot{q})\dot{q} + g(q) + f(\dot{q})$
\end{center}
\begin{align*}
q_{min} &\leq q \leq q_{max}\\
-\dot{q}_{max} &\leq \dot{q} \leq \dot{q}_{max}\\
-\tau_{max} &\leq \tau \leq \tau_{max}\\
\end{align*}
where $q$ is a $7\times1$ vector representing the robot joint angles, $\dot{q}$ is the joint velocity $7\times1$ vector, $\tau$ is the $7\times1$ vector defining joint torques, $M$ is a positive-definite, symmetric $7\times7$ joint inertia matrix, $C$ is a $7\times7$ matrix that represents the Coriolis and centrifugal effects of robot motion, $g$ is a $7\times1$ vector that describes the amount of torque each motor needs to apply to oppose gravity, and $f$ is a $7\times1$ vector that denotes the torques needed to overcome joint friction.

To generate the \textsl{kinodynamic manifold}, a valid end--effector pose is sampled denoting the start configuration of the robot. Random end--effector velocity and duration is also sampled (\Cref{algo:generate-kinodynamic-manifold}, L3-6). These initial states are unrolled forward in time using the manipulator dynamics model while checking for joint position, velocity, acceleration, and torque limits, in addition to workspace limits and collisions with environmental objects (\Cref{algo:generate-kinodynamic-manifold}, L10-14). We additionally introduce a manipulability threshold and nullspace joint velocity sampling to ensure the rollout does not run into singularities or joint velocity limits (\Cref{algo:generate-kinodynamic-manifold}, L17).

% \begin{definition}[Control Projection Function]\label{def:projection}
The \textsl{control projection function} $proj_{MP}$ projects a robot kinodynamic rollout $e \in \mathcal{K}_{MP}$ to a weighted, primitive--specific kinodynamic rollout, i.e.,

\begin{center}
$e^* = proj_{MP}(e)$
\end{center}

where $e^* \in \mathcal{M}_{MP}$.
\end{definition}


Adapting this robot kinodynamic manifold to specific motion primitives requires the introduction of a projection function $\mathcal{proj_{MP}}$ (\Cref{def:projection}). This function leverages a taxonomy of robot motions and environmental characteristics consisting of four interactions:
\begin{itemize}
\item robot only: This considers features of the robot kinodynamic manifold $\mathcal{K}_{MP}$ such as manipulability, dynamic manipulability, and reachability and assigns a weighting on on their quantitative measurements.
\item robot--object: This type of interaction considers the compatibility of robot and object characteristics such as relative gripper--object size and how object height dictates end--effector positioning for nonprehensile manipulation.
\item robot--environment: Here, we ensure that the robot does not collide with movable or fixed obstacles unintentionally. Distribution of clutter and ensuring $e \in \mathcal{K}_{MP}$ has sufficient clearance to execute successfully also informs this metric.
\item object--environment: For nonprehensile interactions, the Coulomb friction between the object and its support surface determines the success of an action. For instance, it is more beneficial for poking to operate in low friction environments to take maximal advantage of its capacity to extend the robot's reachable space. However, the threshold is higher for pushing since constant contact between the robot and the object ensures continuous object motion.
\end{itemize}

Combined, these interactions allow us to quantify the effectiveness of $e \in \mathcal{K}_{MP}$ for a given primitive and operational scenario via $proj_{MP}$, further allowing us to create a \textsl{situated kinodynamic map} for each primitive, grounded in object characteristics and environment configuration.
\ap{more?, see utility.py, task based?, talk about culling edges for poking and grasp action sampling (in planner?)}

% i) robot only, ii) robot--object, iii) robot--environment, iv) object--environment, and, v) task--based.
% \ap{talk about it here, adapt skill + utility, talk about contact taxonomy, what features for each skill, scene--dep weighting}

% \begin{definition}[Situated Kinodynamic Manifold]\label{def:situated-manifold}
A \textsl{situated kinodynamic manifold} is a primitive--specific kinodynamic manifold, grounded in environmental configuration via a projection function $proj_{MP}$
\begin{center}
$\mathcal{M}_{MP} = \{(\mathcal{K}_{MP}, \rho) \mid proj_{MP}(e) \in \mathbb{R}\}$
\end{center}
where $e \in \mathcal{K}_{MP}$.
\end{definition}

% \subsubsection{Real--World Considerations}
% \ap{talk about ramp up and ramp down}

